\clearpage\SourceLeaf{151}%
\pdfbookmark[1]{Gauss an von Lindenau}{lindenau1808aug}
\fancyhead[L]{\small GAUSS AN VON LINDENAU}
\begin{center}
{\small GAUSS AN VON LINDENAU.\par}
\medskip\hrule\medskip
{\small Monatliche Correspondenz zur Beförderung der Erd- und Himmels-Kunde,\par
herausgegeben vom Freiherrn \textsc{von Zach}. September 1808.\par}
\medskip\hrule
\end{center}\medskip
\begin{flushright}Göttingen am 30. August 1808.\end{flushright}
\begin{center}\small
—\ —\ —\ —\ —\ —\ —\ —\ —\ —\ —\ —\ —\ —\ —\ —\ —
\end{center}
Ich nehme mir die Ehre, Ihnen hier für die \emph{Mon. Corresp.} [S.\ 141 d.\ B.] einen Aufsatz über eine Aufgabe zu schicken, deren practische Anwendung sehr empfohlen zu werden verdient. Es scheint mir, dass dieses Verfahren, die Pol-Höhe zu bestimmen, in der Ausübung eine sehr grosse Schärfe verträgt, und wenn man nicht mit Vervielfältigungs-Kreisen zu beobachten Gelegenheit hat, die zuverlässigsten Resultate gibt. Wenigstens fallen alle diejenigen Umstände weg, die bei dem gewöhnlichen Verfahren durch Sonnen-Höhen die Resultate zuweilen zweifelhaft machen können, Ablesen, Fehler des Sextanten, Blendgläser, Glasdach, Refraction. Zwar gibt jede Combination von drei Sternen immer nur ein Resultat, aber wenn man bei Meridian-Höhen der Sonne auch die Beobachtungen noch so sehr vervielfältigt, so bleiben doch die von den erwähnten Umständen herrührenden Fehler, den ersten ausgenommen, immer in ihrer ganzen Stärke zurück; auch hindert ja nichts, dieselben Beobachtungen in mehrern Nächten zu wiederholen und selbst in einer Nacht mehrere Combinationen zu machen, wenn man sie so auswählt, dass die drei Beobachtungen in nicht zu grossen Zwischenräumen auf einander folgen. In der gestrigen Nacht habe ich die drei Sterne, die im Reispiel der Abhandlung angeführt sind, noch einmal in
\SourceLeaf{152}%
\expandafter\def\csname HistoricalNoteText7\endcsname{Fehler des Sextanten bei $105^0\,52'= -49''4$.}%
einer etwas verschiedenen Höhe beobachtet und für die Polhöhe $51^0\,31'\,54''4$\HistoricalNote{7} gefunden. Die drei Resultate sind folgende:
\begin{center}
\begin{tabular}{lr@{\quad}l}
August 25 & $.$ & $51^0\quad31'\quad56''7$\\
27 & $.$ & $51\quad31\quad51.5$\\
30 & $.$ & $51\quad31\quad54.4$\\\cline{1-3}
im Mittel & $.$ & $51\quad31\quad54.2$
\end{tabular}
\end{center}
genau wie \textsc{Mayer} unsere Polhöhe bestimmt hat. Wenn man auch diesen Grad von Übereinstimmung zum Theil hier zufällig halten muss, so glaube ich doch, dass man durch öftere Wiederholung und Vervielfältigung der Beobachtungen mit gut bestimmten Sternen sich der wahren Pol-Höhe immer auf wenige Secunden nähern kann. Alles hängt blos von der Vergrösserung und Deutlichkeit des Fernrohrs und der Achtsamkeit des Beobachters ab. Aus Meridian-Höhen der Sonne, die immer unter sich vortrefflich stimmten, hatte ich zur Zeit der Sonnenwende mehreremal mit einem Sextanten $51^0\,32'\,32''$ gefunden, also über eine halbe Minute fehlerhaft. Ob der Fehler meines Sextanten blos daher rührt, dass der ganze Bogen zu gross ist, habe ich noch nicht hinlänglich geprüft; indess wird das gerade durch die von mir empfohlene Methode sehr leicht künftig geschehen können.

